\chapter{Formula registry}
\begin{longtable}{p{0.29\linewidth}p{0.61\linewidth}}
\toprule
Object & Formula and required interpretation\\
\midrule
Rational Hodge space & $\Hdg^p(X)=H^{2p}(X,\Q)\cap H^{p,p}(X)$.\\
Algebraic space & $\Alg^p(X)=\im(\CH^p(X)_{\Q}\to H^{2p}(X,\Q))$.\\
Real obstruction & $\mathcal O=P_{\Hdg\otimes\R}-P_{\Alg\otimes\R}$. No algebraic realization of these orthogonal projectors is presumed.\\
Cycle frame & $S=CC^*$. Positivity spans a target packet when every supplied algebraic vector lies in it.\\
Unconditional rank test & $\rank_{\Q}[C\ H_{\Q}]=\rank_{\Q}C$ tests target containment in the actual cycle span.\\
Polyphase assembly & $S_{M,r}=\sum_{\ell=0}^{L-1}S_{LM,r+\ell M}$. These are frame sums, not sums of span projectors.\\
Support test & $j^*\alpha=0$. Purity turns it into cycle algebraicity when the support has codimension at least $p$.\\
Divisor transfer & $T=\sum_rB_rA_r$ with actual Chow correspondences, packet preservation and invertibility.\\
Positive transfer & $T=\sum_rA_r^\dagger A_r$ only with algebraic realizations of those adjoints.\\
Balanced Fermat tuple & $\sum_i\langle ua_i\rangle_m=3m$ for every unit $u$.\\
Conductor & $c(a)=m/\gcd(m,a_0,\ldots,a_5)$.\\
Matrix capacity & $C(a)=\sup\{\sum w_v:\sum w_vv=\mathbf1,\ w_v\ge0\}$.\\
Global dual & $y\cdot v\ge1$ for every invariant column implies $C(a)\le y\cdot\mathbf1$.\\
Five-phase surface & $e_1=e_2=P-\kappa z^5=0$, with $\kappa^q=-1/5$.\\
Weil projection & $P_W(c)=-(\langle c,u\rangle u+\langle c,v\rangle v)/32$.\\
\bottomrule
\end{longtable}

\chapter{Correspondence degrees}
Let $X$ have dimension $n$ and $Y$ dimension $m$. A correspondence $\Gamma\in\CH^r(X\times Y)$ acts by
\[
\Gamma_*:H^i(X)\longrightarrow H^{i+2(r-n)}(Y)
\]
and shifts Hodge type by $(r-n,r-n)$. To lower $H^{2p}(X)$ to $H^2(Y)$, take $r=n-p+1$. To return degree two to degree $2p$ on $X$, a correspondence on $Y\times X$ has codimension $m+p-1$.

Transpose incidence alone can have the wrong degree. Insertions of ample classes must be stated, and their effect on the polarization pairing checked. The fact that a cohomological map has a positive adjoint does not identify the adjoint with a cycle correspondence.

For a finite generically finite morphism of degree $d$, the push-pull identity on rational cohomology is $f_*f^*=d\id$. Division by $d$ transfers algebraicity downward. Pullback transfers it upward without division. Coordinate-power maps preserve the old conductor of an inflated character, even though the ambient Fermat degree grows.

\chapter{Source families and chronology}
\begin{longtable}{p{0.29\linewidth}p{0.24\linewidth}p{0.37\linewidth}}
\toprule
Source family & Role & Authority boundary\\
\midrule
August 6 projector/dyadic source & Full infrastructure, Fermat equations and proposed tree argument & Historical proof labels reconciled with the later status note.\\
August 15 Fermat/CM edition & Later scope and construction status & Uniform and correspondence-level write-up obligations retained.\\
September carrier70 v0.1.0 & Normal weighted hypersurface and torus carrier & Exact certificate plus manuscript geometric proposition.\\
September carrier transfer v0.2.0 & Signed degree-210 identity, global dual and character powers & Exact replay and checked classical operations, inheriting carrier geometry.\\
September dual/cycles v0.2.0 & Capacity certificates and explicit standard surfaces & Known Aoki family recomputed, with checked residue theorem.\\
September 22 Weil block & Integral classes, algebraic projector, seed graph and failed deformation route & Exact exterior and tangent calculation, no nonalgebraicity.\\
Preserved Lean v4 & Arithmetic packet formulas and abstract closure theorem & Source inspection, assumed geometric predicate, current build blocked.\\
October review/replay & Corrected reader, exact bounded replays and source contracts & No universal inference from bounded output.\\
\bottomrule
\end{longtable}

The evidence bundle preserves original archives and original source files, keeps current replay outputs separate, and contains the proof audit, computation audit, formal status and checked source contracts. The final reader is a reconciled edition, rather than a silent overwrite of any historical release.

\chapter{How to read the references}
The newly checked theorem leaves include Deligne's rational Hodge statement, the lowest-weight image and semisimplicity results in Hodge II, the uniruled-fourfold input as reproved by Arapura, the BDPP uniruledness criterion, the original Aoki operations and standard equations, the Villaflor complete-intersection period theorem, the stated van Geemen Hodge half-twist construction and the odd-character nonvanishing leaf. Each is used within its actual hypotheses.

The bibliography also preserves broad background references inherited from the August source. Inclusion does not imply that every theorem in those books or papers was freshly extracted and re-audited. The cubic and Gushel--Mukai chapters are classical validation examples with their stated incidence hypotheses. The asymptotic analytic discussion is an explanatory source-dependent branch and does not provide a universal obstruction-vanishing theorem.

Kang's metadata and its relevant attribution in checked later literature are authenticated, while the original full proof was not extracted. The current preprints used for Fermat context were not converted into verified census results. The cited Brosnan implication remains conditional at the exact source-access boundary described above. The local secant-preprint eigencharacter failure is a specific calculation, rather than a global assessment of that paper.

Internal Komorebi seed-calculus references describe the research vocabulary and chronology. They are not substitutes for the actual geometric premises used by this reader's conditional constructions.
